\section{经典线性回归把系数估计接到抽样分布}
\label{sec:11-Mathematical-Statistics/08-Regression-and-ANOVA/01}

\subsection{一条线拟合完了，然后呢}

你记录了一个月里每天空调运行时长 \(x_i\) 和用电量 \(y_i\)。最小二乘给你一条线：
\[
y_i = \beta_0 + \beta_1 x_i + \varepsilon_i.
\]
\(\widehat\beta_0\) 和 \(\widehat\beta_1\) 都算出来了，残差图看起来也没什么异常。现在有人问你：多开一小时空调到底多耗多少电？

你会说 \(\widehat\beta_1\approx 1.2\) 度。对方接着问：你有多大把握它不是 1.0？或者更狠——你确定这不是因为天热所以既开了空调又开了其他电器，而空调本身的贡献其实很小？

第一个问题要的是不确定性量化——斜率估计值周围的误差条。第二个问题要的是因果——系数是否真在测量空调的边际效应，而不只是相关性。

最小二乘给你一个点估计，一个最优投影。但它没告诉你换了另一批日子，斜率会波动多少。要回答这个问题，你需要对数据的生成方式做一个概率假设。

\subsection{从投影到抽样分布：给 \(\varepsilon\) 穿上一套概率外衣}

一般地，回归写成
\[
y = X\beta + \varepsilon, \qquad X\in\mathbb{R}^{n\times p}.
\]
如果 \(X\) 满列秩，最小二乘的解是纯线性代数：
\[
\widehat\beta = (X^\trans X)^{-1}X^\trans y,
\qquad
\widehat y = X\widehat\beta,
\qquad
e = y - \widehat y.
\]
此时已经可以断言 \(e \perp \operatorname{Col}(X)\)——残差与所有列正交，\(\widehat y\) 是 \(y\) 在列空间上的正交投影。这是几何，与概率无关。

现在迈出关键一步：给 \(\varepsilon\) 一个分布。经典假设是
\[
\varepsilon \sim \mathcal{N}(0, \sigma^2 I_n),
\]
即在给定设计矩阵 \(X\) 的条件下，各次观测的误差独立、同方差、零均值、正态。

把 \(y = X\beta + \varepsilon\) 代入 \(\widehat\beta\) 的公式：
\[
\widehat\beta
= \beta + (X^\trans X)^{-1}X^\trans\varepsilon.
\]
\(\widehat\beta\) 是真实 \(\beta\) 加上一个由随机误差驱动的偏移项。因为 \(\varepsilon\) 是正态的，\(\widehat\beta\) 作为它的线性变换也是正态的：
\[
\widehat\beta \sim \mathcal{N}\!\left(\beta,\; \sigma^2 (X^\trans X)^{-1}\right).
\]

这个分布的协方差矩阵由两部分相乘决定：噪声尺度 \(\sigma^2\)，以及 \((X^\trans X)^{-1}\)。前者是数据的内在噪声大小；后者是设计矩阵的几何。如果两列几乎共线，\((X^\trans X)^{-1}\) 的对应对角元会很大——即使 \(\sigma^2\) 不大，系数估计的方差也会被放大。样本多没有用，只要某两个特征高度相关，它们各自系数的单独推断就不可靠。

残差平方和
\[
\operatorname{RSS} = \|y - X\widehat\beta\|_2^2
\]
扣除了 \(p\) 个拟合自由度（最小二乘用了 \(p\) 个参数去追赶数据），因此无偏的噪声方差估计为
\[
s^2 = \frac{\operatorname{RSS}}{n-p}.
\]

现在可以构建第 \(j\) 个系数的检验了。记 \(v_j = [(X^\trans X)^{-1}]_{jj}\)。在 \(H_0: \beta_j = \beta_{j,0}\) 成立且上述所有假设都满足时，
\[
T = \frac{\widehat\beta_j - \beta_{j,0}}{s\sqrt{v_j}} \sim t_{\,n-p}.
\]
分母 \(s\sqrt{v_j}\) 是系数的标准误——它把噪声估计和设计几何合成一个数，告诉你估计值在重复抽样中会波动多少。

相应的置信区间是 \(\widehat\beta_j \pm t_{n-p,\,1-\alpha/2}\cdot s\sqrt{v_j}\)。区间宽度随 \(n\) 增大而收缩（\(t\) 分位数趋向正态分位数，\(s\) 更精确），也随共线性加剧而膨胀。

\subsection{\(R^2\) 衡量的是分解比例，不是真理比例}

模型含截距时，总平方和可以按投影几何分解：
\[
\operatorname{TSS} = \sum_i (y_i - \bar y)^2,
\qquad
\operatorname{RSS} = \sum_i (y_i - \widehat y_i)^2,
\]
\[
\operatorname{TSS}
= \underbrace{\sum_i (\widehat y_i - \bar y)^2}_{\operatorname{ESS}} + \operatorname{RSS},
\qquad
R^2 = 1 - \frac{\operatorname{RSS}}{\operatorname{TSS}}.
\]

\(R^2\) 说的是：在这批样本、这个尺度、这些特征下，响应变量的变化中有多大比例被回归平面捕获了。仅此而已。

它不说的东西有一长串：未来数据上的预测精度；残差是否满足模型假设；变量之间的关系是否有因果含义；是不是漏掉了关键的混杂因素。往模型里塞任何变量——哪怕随机生成的噪声列——训练 \(R^2\) 都只升不降。因为增加列只扩大了可投影子空间，残差范数不可能变大。

调整 \(R^2\) 对增加变量做了惩罚，但它仍然是样本内指标。交叉验证和独立测试集评估的是样本外预测能力。三种工具各管一摊，不能互相冒充。

\subsection{嵌套模型的 \(F\) 检验：新增方向是否抓住了真信号}

设约束模型有 \(p_0\) 个参数，完整模型有 \(p_1 > p_0\) 个参数，残差平方和分别是 \(\operatorname{RSS}_0\) 和 \(\operatorname{RSS}_1\)。因为完整模型的列空间包含了约束模型的列空间，投影的子空间更大，\(\operatorname{RSS}_1 \le \operatorname{RSS}_0\) 是几何必然。

但这个下降是真实的信号捕获，还是纯粹因为多加了几个自由度、拟合了部分噪声？\(F\) 统计量把下降量除以新增参数个数、再除以噪声尺度：
\[
F = \frac{(\operatorname{RSS}_0 - \operatorname{RSS}_1) / (p_1 - p_0)}
{\operatorname{RSS}_1 / (n - p_1)}.
\]

分子是``每个新增参数平均带来的平方和下降''，分母是完整模型中的噪声水平。只有当下降显著大于背景噪声波动时，才拒绝约束模型。在零假设（新增方向真实系数全为零）和全部经典假设下，\(F \sim F_{p_1-p_0,\,n-p_1}\)。

单个系数的 \(t\) 检验是这里的特例：此时 \(F = T^2\)，两个检验等价。

\subsection{哪些条件可以松，哪些结论会跟着变}

逐条检查假设及其被违反后的后果，比记住公式更有用。

零条件均值 \(\mathbb{E}[\varepsilon\mid X] = 0\) 是系数估计无偏性的支柱。违反而来的不是方差变大，而是估计在朝一个错误的目标收敛——这叫遗漏变量偏差，增大样本量不能修复。

同方差和误差不相关是标准误公式的前提。异方差时，标准误公式给出的是错误宽度——真实覆盖概率偏离名义水平。稳健协方差估计（sandwich estimator）可以修正标准误，但不能修正系数估计本身的偏移。时间相关或聚类结构需要用相应的相关结构标准误或混合效应模型，否则标准误向乐观方向偏。

Gaussian 假设给了你小样本下精确的 \(t\) 和 \(F\) 分布。大样本下中心极限定理出手——即使误差非正态，\(t\) 统计量也趋向正态。但小样本加厚尾误差时，t 检验的尾部概率可能严重失准。

最难察觉也最常见的问题：模型是在看过数据后筛选出来的。如果你尝试了 30 个变量组合，选了 \(R^2\) 最高的那个再报告其系数的 \(p\) 值和区间——这些区间和 \(p\) 值没有把筛选过程纳入分布计算。它们假定模型是预先指定的，而非从数据中挑出来的。同样的数据先用来选模型、再用来推断模型，推断的分布不再是纸面上的 \(t\) 和 \(F\)。

如果目标是预测：把完整的特征工程、模型选择和调参流程打包放进交叉验证循环，让每一个决策都在训练折叠内完成，验证折叠只用于评估。如果目标是解释系数：在研究设计阶段就声明模型、报告效应量和区间而非只报显著性，检查残差图、杠杆点和影响点，尝试合理的替代规格（加入可能的混杂、变换函数形式）看结论是否稳健。回归推断评价的是整套建模程序的质量，不是最后一次矩阵运算的数值结果。
